\documentclass[10pt,a4paper]{article}

% Черновик для правки названий. Источник — src/content/sections.ts.

% ---------- шрифты и язык ----------
\usepackage{cmap}
\usepackage[T2A]{fontenc}
\usepackage[utf8]{inputenc}
\usepackage[russian]{babel}
\usepackage{amsmath,amssymb}
\usepackage{tempora}

% ---------- поля ----------
\usepackage[a4paper,left=14mm,right=14mm,top=14mm,bottom=14mm]{geometry}

% ---------- таблицы ----------
\usepackage{array,booktabs,longtable}
\usepackage{ragged2e}
\usepackage[table]{xcolor}

\definecolor{rule}{HTML}{C9CDC7}
\arrayrulecolor{rule}
\renewcommand{\arraystretch}{1.5}
\setlength{\parindent}{0pt}
\flushbottom

\begin{document}

{\LARGE\bfseries Молекулярная физика}\par
\vspace{2mm}
{\large Названия карточек}\par
\vspace{3mm}
{\small Формулы перенесены из шпаргалки, названия — из приложения.
Правьте второй столбец; формулы тоже можно менять.}\par
\vspace{6mm}

\begin{longtable}{@{}>{\RaggedRight\arraybackslash}p{0.46\linewidth}@{\hspace{8mm}}>{\RaggedRight\arraybackslash}p{0.46\linewidth}@{}}
\toprule
\textbf{Формула} & \textbf{Название карточки} \\
\midrule
\endfirsthead
\toprule
\textbf{Формула} & \textbf{Название карточки} \\
\midrule
\endhead
\bottomrule
\endfoot
$\displaystyle \nu=\dfrac{N}{N_A}$ & Количество вещества (через число частиц) \\
\addlinespace[2mm]
$\displaystyle M=m_0 N_A$ & Молярная масса (через массу молекулы) \\
\addlinespace[2mm]
$\displaystyle m=m_0 N$ & {Масса газа (через массу молекулы)} \\
\addlinespace[2mm]
$\displaystyle n=\dfrac{N}{V}$ & Концентрация молекул газа \\
\addlinespace[2mm]
$\displaystyle \nu=\dfrac{m}{M}$ & Количество вещества (через массу тела) \\
\addlinespace[2mm]
$\displaystyle m_0=\dfrac{m}{N}=\dfrac{M}{N_A}$ & Масса одной молекулы \\
\addlinespace[2mm]
$\displaystyle N=\nu N_A=\dfrac{m}{M}N_A$ & Число молекул \\
\addlinespace[2mm]
$\displaystyle n=\dfrac{N}{V}=\dfrac{\nu N_A}{V}$ & Концентрация молекул (через количество вещества) \\
\addlinespace[2mm]
$\displaystyle \rho=\dfrac{m}{V}=m_0 n$ & Плотность газа (через концентрацию молекул) \\
\midrule
\addlinespace[2mm]
$\displaystyle p=\dfrac{1}{3}n m_0\langle v^2\rangle$ & Основное уравнение МКТ \\
\addlinespace[2mm]
$\displaystyle \langle E_{\text{к}}\rangle=\dfrac{m_0\langle v^2\rangle}{2}$ & Средняя кинетическая энергия поступательного движения молекулы \\
\addlinespace[2mm]
$\displaystyle p=\dfrac{2}{3}n \langle E_{\text{к}}\rangle$ & Давление газа через среднюю кинетическую энергию молекулы \\
\addlinespace[2mm]
$\displaystyle T=t+273$ & Перевод температуры из шкалы Цельсия в Кельвина \\
\addlinespace[2mm]
$\displaystyle p=nkT$ & Давление газа через концентрацию и температуру \\
\addlinespace[2mm]
$\displaystyle \langle E_{\text{к}}\rangle=\dfrac{3}{2}kT$ & Средняя кинетическая энергия молекулы через температуру \\
\addlinespace[2mm]
$\displaystyle R =k N_A = 8,31 $ Дж/(моль$\cdot$K) & Универсальная газовая постоянная \\
\addlinespace[2mm]
$\displaystyle v_{\text{кв}}=\sqrt{\langle v^2\rangle}$ & Средняя квадратичная скорость \\
\addlinespace[2mm]
$\displaystyle v_{\text{кв}}=\sqrt{\dfrac{3kT}{m_0}}$ & Средняя квадратичная скорость (через температуру) \\
\addlinespace[2mm]
$\displaystyle v_{\text{кв}}=\sqrt{\dfrac{3RT}{M}}$ & Средняя квадратичная скорость (через молярную массу) \\
\midrule
\addlinespace[2mm]
$\displaystyle pV=\dfrac{m}{M}RT=\nu RT$ & Уравнение Клапейрона—Менделеева \\
\addlinespace[2mm]
$\displaystyle \dfrac{pV}{T}=\text{const}$ & Уравнение Клапейрона \\
\addlinespace[2mm]
$\displaystyle \dfrac{p_1V_1}{T_1}=\dfrac{p_2V_2}{T_2}$ & Уравнение Клапейрона для двух состояний \\
\addlinespace[2mm]
$\displaystyle p_{\text{смеси}}=p_1+p_2+\dots$ & Закон Дальтона \\
\addlinespace[2mm]
$\displaystyle p=\dfrac{\rho RT}{M}$ & Давление через плотность \\
\addlinespace[2mm]
$\displaystyle pV=\text{const}$ & Закон Бойля—Мариотта \\
\addlinespace[2mm]
$\displaystyle \dfrac{V}{T}=\text{const}$ & Закон Гей-Люссака \\
\addlinespace[2mm]
$\displaystyle \dfrac{p}{T}=\text{const}$ & Закон Шарля \\
\midrule
\addlinespace[2mm]
$\displaystyle U=\dfrac{3}{2}\dfrac{m}{M}RT=\dfrac{3}{2}pV$ & Внутренняя энергия идеального газа \\
\addlinespace[2mm]
$\displaystyle \Delta U=\dfrac{3}{2}\nu R\Delta T$ & Изменение внутренней энергии \\
\addlinespace[2mm]
$\displaystyle A=p\Delta V=p(V_2-V_1)$ & Работа газа при постоянном давлении \\
\addlinespace[2mm]
$\displaystyle Q=cm\Delta T=cm(T_2-T_1)$ & Нагревание и охлаждение \\
\addlinespace[2mm]
$\displaystyle Q=\lambda m$ & Плавление и кристаллизация \\
\addlinespace[2mm]
$\displaystyle Q=Lm$ & Парообразование и конденсация \\
\addlinespace[2mm]
$\displaystyle Q=qm$ & Сгорание топлива \\
\addlinespace[2mm]
$\displaystyle Q=\Delta U+A=\Delta U+p\Delta V$ & Первый закон термодинамики \\
\addlinespace[2mm]
$\displaystyle \eta=\dfrac{A_{\text{п}}}{Q_1}=\dfrac{Q_1-Q_2}{Q_1}$ & КПД теплового двигателя \\
\addlinespace[2mm]
$\displaystyle Q_1=Q_2+A_{\text{п}}$ & Баланс энергии за цикл \\
\addlinespace[2mm]
$\displaystyle \varphi=\dfrac{p_{\text{п}}}{p_{\text{н}}}\cdot100\%$ & Относительная влажность через давление \\
\addlinespace[2mm]
$\displaystyle \varphi=\dfrac{\rho_{\text{п}}}{\rho_{\text{н}}}\cdot100\%$ & Относительная влажность через плотность \\
\addlinespace[2mm]
\end{longtable}

\vspace{4mm}
{\footnotesize Всего карточек: 39.}

\end{document}
